\section{Analytic foundations and conventions}
\label{sec:r10-foundations}
The Fourier transform is unitary, $D=-i\nabla$, and
$P(u)=\operatorname{Im}\int\overline u\nabla u=\int\xi|\widehat u|^2$.
The real Hilbert pairing is the real part of the complex $L^2$ pairing.
Throughout the local-power arguments, $N\ge1$, $m>1$, $s=m/2$,
$r=m/N$, and $p=2+2r$. A principal symbol $a$ is real, even,
$m$-homogeneous, $C^{N+2}$ off zero, and satisfies
$c|\xi|^m\le a(\xi)\le C|\xi|^m$.
Set $A(u)=\int a|\widehat u|^2$, $M(u)=\|u\|_2^2$, and
$E(u)=A(u)/2-\|u\|_p^p/p$. The sharp critical mass $M_c$ is defined by
\begin{equation}\label{eq:r10-basic-GN}
 \|u\|_p^p\le\frac p2\left(\frac{M(u)}{M_c}\right)^r A(u).
\end{equation}
The traveling-wave convention $u(t,x)=e^{i\omega t}U(x-vt)$ gives
$[a(D)+\omega+iv\cdot\nabla]U-|U|^{2r}U=0$.

\begin{lemma}[Normalized optimizer compactness]
The inequality \eqref{eq:r10-basic-GN} is attained. If
$M(f_n)\to M_c$, $A(f_n)\to1$, and $\|f_n\|_p^p\to p/2$, then
some translated subsequence converges strongly in $H^s$ to an optimizer.
The normalized optimizer set is compact modulo translations and phases.
\end{lemma}
\begin{proof}
Sharp Sobolev interpolation gives a finite positive best constant, and
amplitude and spatial scaling normalize a maximizing sequence as stated.
Bounded energy-space sequences that vanish locally in mass converge to
zero in $L^p$: use a bounded-overlap cover by unit balls, local Sobolev
interpolation, and $2<p<2N/(N-m)$ when $N>m$ (all finite $p$ otherwise).
Thus translations give a nonzero weak limit $Q$ and almost-everywhere
convergence. Write $M_1=M(Q)>0$, $A_1=A(Q)>0$, $M_2=M_c-M_1$,
and $A_2=1-A_1$. Quadratic splitting and Brezis--Lieb give
\[
 M_c^r\le M_1^r A_1+M_2^r A_2,\qquad A_1+A_2=1.
\]
If both masses are positive, the right side is strictly below $M_c^r$.
Consequently $M_1=M_c$ and then $A_1=1$. Weak convergence plus
convergence of both quadratic norms gives strong $H^s$ convergence.
The same proof applies to a sequence of optimizers. Constant phase is
proved by the density-norm lemma below; this compactness argument itself
does not need a positivity-preserving kinetic operator.
\end{proof}

\begin{lemma}[Weak-profile regularity and Pohozaev identities]
Every weak $H^s$ local-power profile belongs to $L^\infty\cap H^{m+1}$
and satisfies
\begin{equation}\label{eq:r10-basic-pohozaev}
 v\cdot P=2mE,\qquad \omega M=rA+(2m-p)E.
\end{equation}
No spatial moment is required.
\end{lemma}
\begin{proof}
Choose $\Lambda$ so large that
$a(\xi)-v\cdot\xi+\Lambda\asymp1+|\xi|^m$.
The shifted resolvent maps $L^q$ to $H^{m,q}$, $1<q<\infty$,
by the Mihlin multiplier theorem. For $N>m$, the starting exponent
is $q_0=2N/(N-m)$. Splitting the nonlinearity into $|U|>1$ and
$|U|\le1$, resolvent Sobolev improvement iterates the reciprocal
exponent as $\beta_{j+1}=(1+2r)\beta_j-r$ with
$\beta_0=(1-r)/2$. The ideal iterates
$\beta_j=1/2-(r/2)(1+2r)^j$ cross the boundedness threshold after
finitely many steps. Inputs have exponent greater than one; intermediate
exponents may be weakened to avoid endpoints. The linear contribution
has the stronger improvement $\beta\mapsto\beta-r$ and is handled
separately. When $N=m$ start at any sufficiently large finite exponent;
when $N<m$ the initial space already embeds in $L^\infty$.
Now the nonlinearity is in $L^2$, giving $U\in H^m\subset H^1$.
The $C^1$ map $z\mapsto|z|^{2r}z$, with bounded derivative on the
range of $U$, sends $U$ into $H^1$; one more resolvent application
gives $H^{m+1}$.

Here are the cutoff details for the dilation calculation. Let
$f_R(x)=x\chi(x/R)$ with $\chi$ compactly supported and equal to
one near zero, and let $L_R=f_R\cdot\nabla+\operatorname{div}f_R/2$.
The Fourier kernel of $[a(D),L_R]$ is, up to the transform constant,
\[
 \frac i2(a(\xi)-a(\eta))(\xi+\eta)\cdot\widehat f_R(\xi-\eta).
\]
Splitting comparable and noncomparable frequencies gives
\[
 |a(\xi)-a(\eta)||\xi+\eta|
 \le C|\xi-\eta|(|\xi|+|\eta|)^m.
\]
After division by $\langle\xi\rangle^s\langle\eta\rangle^s$,
the convolution kernel is bounded by
$C|k|\langle k\rangle^s R^{N+1}|\widehat f(Rk)|$, with a
uniform $L^1$ norm for $R\ge1$. Hence the commutators are uniformly
bounded $H^s\to H^{-s}$. Schwartz approximation and the homogeneous
symbol identity give
$\langle U,[a(D),L_R]U\rangle\to mA(U)$.
The drift commutator tends to $-v\cdot P$; its order is one and
$s>1/2$. In the nonlinear pairing integration by parts and dominated
convergence give
\[
 \operatorname{Re}\int |U|^{p-2}\overline U L_RU
 \longrightarrow N(1/2-1/p)\|U\|_p^p=\frac m p\|U\|_p^p.
\]
Testing the profile with $L_RU$ therefore gives $mE=v\cdot P/2$.
Testing with $U$ gives $A+\omega M-v\cdot P-\|U\|_p^p=0$;
combining these two equations proves \eqref{eq:r10-basic-pohozaev}.
\end{proof}

The same regularity proof holds for an individually fixed nonnegative
order-$m$ addition $b$ with the stated Mihlin derivative bounds.
Constants need not be uniform in that addition or in $v$.
For nonhomogeneous minimizers, differentiation of their scalar
mass-preserving dilation energy is legitimate directly by dominated
convergence and minimality; differentiability of the dilation orbit
as an $H^s$-valued map is not being assumed.
